\begin{remarks}[6.8.4]
\label{IV.2.s6.8.4-fr}
\begin{enumerate}
\item[(i)] Si $f$ es fielmente plano, $g$ es plano y si $g\circ f$ es de
coprofundidad $\leq n$, entonces $g$ es de coprofundidad $\leq n$, como
resulta de (6.6.2).
\item[(ii)] Cuando, en (6.8.1), se toma para $Y$ el espectro de un cuerpo
$k$, las nociones (iv), (v) y (vi) se reducen a las definidas en (6.7.5).
Es claro que estas últimas son relativas al cuerpo de base $k$. La
definición (6.8.1) lleva entonces a decir que un preesquema $X$ es «\emph{regular}
(resp. $(R_n)$, resp. \emph{normal}) \emph{sobre} $k$» en lugar de decir que es
«geométricamente regular (resp. $(R_n)$, resp. normal) con respecto a $k$»;
se evitará confundir esta noción con la propiedad de ser regular (resp.
$(R_n)$, resp. normal), que es independiente de $k$. A este respecto se
pueden hacer las mismas observaciones que en (4.5.12).
\end{enumerate}
\end{remarks}

\begin{proposition}[6.8.5]
\label{IV.2.s6.8.5-fr}
Sean $k$ un cuerpo, $X$, $Y$ dos $k$-preesquemas localmente noetherianos,
uno de los cuales es localmente de tipo finito sobre $k$. Para un $k$-preesquema $Z$,
designemos por $\boldsymbol{P}(Z)$ una de las propiedades \textup{c)} a \textup{f)} de \textup{(6.6.1)}; la propiedad
«$\boldsymbol{P}(Z)$ geométrica» se define entonces en \textup{(6.7.6) (resp. (4.6.1))} cuando
$\boldsymbol{P}(Z)$ es una de las propiedades \textup{c), d), e) (resp. f))} de \textup{(6.6.1)}. Entonces:
\begin{enumerate}
\item[\textup{(i)}] Si $X$ posee la propiedad $\boldsymbol{P}$, e $Y$ la propiedad $\boldsymbol{P}$
geométrica, $X\times_kY$ posee la propiedad $\boldsymbol{P}$.
\item[\textup{(ii)}] Si $X$ e $Y$ poseen la propiedad $\boldsymbol{P}$ geométrica, también la
posee $X\times_kY$.
\end{enumerate}
\end{proposition}

\begin{proof}
Sean, en efecto, $f:X\times_kY\to X$, $g:X\to\operatorname{Spec}(k)$ los
morfismos estructurales, que son fielmente planos (2.2.13). La hipótesis de que
$Y$ posee la propiedad $\boldsymbol{P}$ geométrica implica, en virtud de (6.8.2), que
se cumple $\boldsymbol{M}(f)$, donde $\boldsymbol{M}$ es la propiedad de (6.8.3) que corresponde a $\boldsymbol{P}$; bajo
la hipótesis (ii), también se cumple $\boldsymbol{M}(g)$, por lo que la afirmación (ii) resulta de
(6.8.3, (i)). En cuanto a la afirmación (i), resulta directamente de (6.6.1, (i)).
\end{proof}
